\section{A Reality Check and Proposed Work}
\elaine{TODO: fix table 1}

The goal of the proposal is to enable private accesses to on-chain data
at the Ethereum scale. 
We first perform a reality check of where we are today 
by comparing several 
alternative techniques
in terms of their performance and security assumptions (\Cref{sec:comparison}), 
and then give an overview our proposed work (\Cref{sec:proposed}). 



% Comparison table, transcribed from table.pdf.
\begin{table}[H]
\centering
\footnotesize
\renewcommand{\arraystretch}{1.1}
\setlength{\tabcolsep}{3pt}
\newcommand{\q}[1]{\linebreak{\scriptsize #1}}% qualifier under a value
\newcommand{\sysrow}[2]{\textbf{#1}\linebreak{\scriptsize #2}}
\newcommand{\band}[1]{%
  \rowcolor{black!8}\multicolumn{5}{l}{\rule{0pt}{2.4ex}\textbf{#1}}\\}
\newcommand{\tn}[1]{\textsuperscript{#1}}
\newcommand{\schemesep}{% gray rule between schemes
  \arrayrulecolor{black!25}\specialrule{0.4pt}{3pt}{3pt}\arrayrulecolor{black}}
\caption{Private access to Ethereum data: existing approaches and the
  targets of our proposed work. Our existing ORAM and Garbled RAM were run on a 32-core Intel Xeon
   Platinum 8352S machine.}
\label{tab:comparison}
\begin{tabularx}{\linewidth}{>{\raggedright\arraybackslash}p{3.4cm}
                             *{4}{>{\centering\arraybackslash}X}}
\toprule
& \multicolumn{2}{c}{\textbf{Ethereum accounts}}
& \multicolumn{2}{c}{\textbf{Full Ethereum state}} \\
& \multicolumn{2}{c}{\scriptsize ${\sim}206.8$M records of 64 bytes}
& \multicolumn{2}{c}{\scriptsize $2^{32}$ records of 64 bytes} \\
\cmidrule(lr){2-3}\cmidrule(lr){4-5}
& Latency & Throughput & Latency & Throughput \\
& & {\scriptsize per GPU or CPU} & & {\scriptsize per GPU or CPU} \\
\specialrule{\lightrulewidth}{\aboverulesep}{0pt}
\band{Status quo}
\addlinespace[3pt]
\sysrow{neOMAP}{Oblivious Labs' TEE-based ORAM~\cite{neomap}}
& 30\,$\mu$s\q{non-batched}\par\smallskip
  10--20\,ms\q{batched}
& 33K\q{queries/sec}\par\smallskip
  300K--400K\q{queries/sec}
& 50\,$\mu$s\q{non-batched}\par\smallskip
  50\,ms\q{batched}
& 20K\q{queries/sec}\par\smallskip
  33K\tn{$\star$}\q{queries/sec} \\
\schemesep
\sysrow{InsPIRe}{Ethereum's GPU-based PIR~\cite{efinspire}}
& ${\sim}250$\,ms
& 40--60\q{queries/sec}
& 5.2\,s\tn{$\dagger$}
& 2--3\q{queries/sec}\tn{$\dagger$} \\
\schemesep
\sysrow{Garbled RAM}{our preliminary work~\cite{picogram,zebra}}
& ${\sim}400$\,ms\tn{$\dagger$}
& 1/latency\tn{$\diamond$}%\tn{$\dagger$}
& ${\sim}700$\,ms\tn{$\dagger$}
& 1/latency\tn{$\diamond$} \\
\addlinespace[3pt]
\band{Proposed work (ballpark targets)}
\addlinespace[3pt]
\sysrow{MicrORAM}{faster TEE-based ORAM (Thrust~2)}
& ${<}2$\,$\mu$s\q{non-batched}\par\smallskip
  ${<}300$\,$\mu$s\q{batched}
& ${>}2$M\q{queries/sec}\q{batched}
& ${<}3$\,$\mu$s\q{non-batched}\par\smallskip
  ${<}400$\,$\mu$s\q{batched}
& ${>}1.5$M\q{queries/sec}\q{batched} \\
\schemesep
\sysrow{Garbled RAM}{our proposed work (Thrust~1)}
& ${\sim}100$\,ms
& 40--60\q{queries/sec}
& ${\sim}150$\,ms
& 30--40\q{queries/sec} \\
\addlinespace[2pt]
\bottomrule
\end{tabularx}

\smallskip
\begin{minipage}{\linewidth}
\scriptsize
\tn{$\star$}\,neOMAP swaps data to SSD at this scale, and the throughput is bounded by the SSD bandwidth. \\
% \tn{$\star$}\,Could be faster with parallelization; the database no longer fits
% in a single GPU.\quad
% \tn{$\diamond$}\,Cannot be parallelized away; the actual number can be worse because
% the database does not fit in a single GPU.\quad
\tn{$\diamond$}\,The current implementation has no throughput optimization. \\
\tn{$\dagger$}\,Extrapolated, current experiments are at smaller scales.
\end{minipage}
\end{table}

\begin{samepage}
\subsection{Reality Check}
\label{sec:comparison}

Three approaches offer different tradeoffs between performance, scalability,
and trust for private access to blockchain data: trusted execution environments
(TEEs) combined with Oblivious RAM (ORAM), private information retrieval (PIR),
and Garbled RAM. The comparison in \Cref{tab:comparison} considers both an account-sized database
and a larger database representing the full Ethereum state. It illustrates why
we propose to advance two complementary directions: faster TEE-based access for
deployments that accept trusted hardware, and scalable cryptographic access for
deployments that require privacy without that hardware assumption.
\Cref{tab:cost} translates the throughput in \Cref{tab:comparison} into the
number of machines, and the annual cost on AWS, needed when 10\% of a
typical RPC load requires privacy.
\par
\end{samepage}

% Cost of serving a typical RPC load, derived from Table 1's throughput.
\begin{table}[H]
\centering
\footnotesize
\renewcommand{\arraystretch}{1.1}
\setlength{\tabcolsep}{3pt}
\newcommand{\q}[1]{\linebreak{\scriptsize #1}}% qualifier under a value
% Row padding sits inside the rows rather than in \addlinespace or rule
% spacing, so the shaded cost columns have no gaps.
\newcommand{\toppad}{\rule{0pt}{\dimexpr\arraystretch\ht\strutbox+3pt\relax}}
\newcommand{\botpad}{\rule[-\dimexpr\arraystretch\dp\strutbox+3pt\relax]{0pt}{0pt}}
\newcommand{\sysrow}[2]{\toppad\textbf{#1}\linebreak{\scriptsize #2}\botpad}
\newcommand{\band}[1]{%
  \rowcolor{black!8}\multicolumn{7}{l}{\rule{0pt}{2.4ex}\textbf{#1}}\\}
\newcommand{\schemesep}{% gray rule between schemes
  \arrayrulecolor{black!25}\specialrule{0.4pt}{0pt}{0pt}\arrayrulecolor{black}}
\newcommand{\cpu}{m6id.\linebreak 16xlarge}
\newcommand{\largecpu}{x8i.\linebreak 24xlarge}
\newcommand{\gpu}{p5.\linebreak 4xlarge}
\captionsetup{justification=raggedright}% long URLs
\caption{Machines and annual AWS cost for 10K private requests/sec
  (10\% of a typical 100K/sec RPC load). Flat-database columns assume one
  lookup per request. The added account-MPT columns include complete
  proof retrieval for 206.8M accounts; they are illustrative extrapolations
  under the model below, not benchmarks. K, M, and B denote thousands,
  millions, and billions.}
\label{tab:cost}
\begin{tabularx}{\linewidth}{>{\raggedright\arraybackslash}p{2.65cm}
                             *{3}{>{\centering\arraybackslash}X
                                  >{\columncolor{yellow!20}\centering\arraybackslash}X}}
\toprule
& \multicolumn{2}{c}{\textbf{Ethereum accounts}}
& \multicolumn{2}{c}{\textbf{Full Ethereum state}}
& \multicolumn{2}{c}{\textbf{Accounts + MPT proof}} \\
& \multicolumn{2}{c}{\scriptsize ${\sim}206.8$M records of 64 bytes}
& \multicolumn{2}{c}{\scriptsize $2^{32}$ records of 64 bytes}
& \multicolumn{2}{c}{\scriptsize ${\sim}206.8$M accounts} \\
\cmidrule(lr){2-3}\cmidrule(lr){4-5}\cmidrule(lr){6-7}
& Config & Annual cost & Config & Annual cost & Config & Annual cost \\[\aboverulesep]
\specialrule{\lightrulewidth}{0pt}{0pt}
\band{Status quo}
\sysrow{neOMAP}{Oblivious Labs' TEE-based ORAM}
& \cpu\q{1 machine}\q{batched}
& \$33K\q{batched}
& \largecpu\q{1 machine}\q{batched}
& \$92K\q{batched}
& \cpu\q{6 capacity units}
& \$200K \\
\schemesep
\sysrow{InsPIRe}{Ethereum's GPU-based PIR}
& \gpu\q{167--250 machines}
& \$10M--\$15M
& \gpu\q{3.3K--5K machines}
& \$201M--\$301M
& \gpu\q{20K--30K machines}
& \$1.21B--\$1.81B \\
\schemesep
\sysrow{Garbled RAM}{our preliminary work}
& \cpu\q{4K machines}
& \$133M
& \cpu\q{7K machines}
& \$233M
& \cpu\q{24K capacity units}
& \$799M \\
\band{Proposed work (ballpark targets)}
\sysrow{MicrORAM}{faster TEE-based ORAM (Thrust~2)}
& \cpu\q{1 machine}\q{batched}
& \$33K\q{batched}
& \largecpu\q{1 machine}\q{batched}
& \$92K\q{batched}
& \cpu\q{6 capacity units}
& \$200K \\
\schemesep
\sysrow{Garbled RAM}{our proposed work (Thrust~1)}
& \cpu\q{167--250 machines}
& \$5.6M--\$8.3M
& \cpu\q{250--334 machines}
& \$8.3M--\$11M
& \cpu\q{1,002--1,500 capacity units}
& \$33.4M--\$49.9M \\
\specialrule{\heavyrulewidth}{0pt}{\belowbottomsep}
\end{tabularx}
\end{table}

\begingroup
\footnotesize
\raggedright
\paragraph{Cost and load assumptions.}
The 100K/sec RPC load is consistent with the average, across-chain rates
reported by major providers: Infura, over 13B requests/day
(${\approx}150$K/sec, \url{https://www.theblock.co/linked/42077/consensys-acquires-ethereum-infrastructure-service-provider-infura});
Alchemy, over 10B/day (${\approx}120$K/sec,
\url{https://us.ovhcloud.com/resources/case-studies/alchemy-scales-blockchain-with-ovhclouds-bare-metal/});
QuickNode, over 8.2B/day in 2024 (${\approx}95$K/sec,
\url{https://www.quicknode.com/blog/quicknodes-2024-wrapped}); and Ankr,
1T/month (${\approx}390$K/sec,
\url{https://www.ankr.com/blog/inside-ankr-s-1-t-monthly-rpc-requests-served-powering-web3/}).
We retain the existing October 2026 us-east-1 planning rates
(\url{https://aws.amazon.com/ec2/pricing/on-demand/}): m6id.16xlarge at
\$3.80/hour (\$33,288/year), p5.4xlarge at \$6.88/hour
(\$60,268.80/year), and x8i.24xlarge at \$10.50/hour
(\$91,980/year), with 8,760 hours of operation annually. These are the
existing CPU and H100 instance choices; the full-state neOMAP and
MicrORAM columns retain the larger x8i instance. Flat-database machine
counts divide the private load by the throughput in \Cref{tab:comparison}
and round up.
\par\endgroup

\paragraph{Account-MPT cost model (illustrative).}
The new columns in \Cref{tab:cost} use only the 206.8M-account case.
We model a resident pruned trie with $N=206.8\times10^6$ retained account
leaves, rooted at Ethereum's state root and retaining the boundary hashes
of omitted sub-tries. All nodes on the requested proof path are retrieved
privately; no client-side proof cache is assumed.
A canonical account proof contains the RLP-encoded nodes on the path,
including the sibling references in each branch~\cite{eip1186,ethmpt}.
Reading just the selected 32-byte child reference suffices for navigation,
but cannot supply that complete proof. The ORAM and Garbled RAM layouts
store node encodings in 32-byte pieces and retrieve each proof piece once;
there is no additional full-node read just to select the next child.

For a representative path, take eight branch nodes, one extension
(two pieces), and one account leaf (five pieces, including its encoded
path and all four account fields). Eight branches is a rounded planning
estimate for hashed keys at this scale, not a maximum-depth guarantee.
A fully populated branch encodes in at most 532 bytes, or 17 pieces.
Thus the proof workload is $k=8\cdot17+2+5=143$ piece reads per request
(4,576 bytes before transport overhead). This budgets full branches;
sparse branches can be smaller. Stored encodings and hashes are reused,
so serving a proof does not require recomputing the trie or evaluating
Keccak inside a garbled circuit. Private deployments must pad paths and
proof lengths to their chosen public limits; paths exceeding this
representative budget require additional work.

We use a storage expansion $s=6$ relative to the flat $64N$-byte database,
including leaf data, child references, and node metadata. This is a
planning budget, not a measured Ethereum footprint. For orientation, a
random hexary trie at this size has about $0.33N$ branches; reserving 17
pieces per branch and five per leaf already costs about $5.3\times64N$
bytes. The remaining allowance covers extensions and layout overhead.
The resulting logical database is about 79.4\,GB, or $12N$ 32-byte pieces.
All references into omitted sub-tries must be retained. The 206.8M-account
baseline is not itself a census of every leaf in Ethereum's state trie;
an actual pruned layout must be measured before deployment.

At $L=10{,}000$ private requests/sec, we assume neOMAP, MicrORAM,
and both Garbled RAM configurations are memory-capacity-bound. Their MPT
costs therefore scale with storage, rather than with the number of proof
reads. If $m_0$ and $C_0$ are the account-case provisioned capacity and cost
in the flat-database columns, the model is
\[
 m_{\mathrm{MPT}}=6m_0, \qquad C_{\mathrm{MPT}}\approx6C_0.
\]
We apply this to the existing rounded-up machine counts, with costs
calculated at the unrounded annual instance rate. This gives six capacity
units and \$199,728/year for each TEE scheme; 24,000 units and
\$798,912,000/year for preliminary Garbled RAM; and 1,002--1,500 units and
\$33,354,576--\$49,932,000/year for proposed Garbled RAM. A capacity unit
is the memory provisioned by one m6id.16xlarge at the retained rate.
These entries assume sixfold memory provisioning remains sufficient for
the proof workload, a constant physical-to-logical storage factor, and
private access across any partitions. They are stipulated capacity
estimates, not new throughput measurements; full replicas alone do not
expand the addressable database. Garbled RAM retains the baseline's
protocol-party accounting. Larger memory instances may change cost.

\emph{InsPIRe with whole-node responses.}
InsPIRe remains compute-bound. Its GPU engine accepts configurable entry
widths and recommends padded widths of 15, 30, 60, 120, 240, \ldots\ bytes;
the serving layer already packs multiple cells into a response
record~\cite{inspiregpu,efinspire}. We propose separate branch, extension,
and leaf databases with whole-node records. For example, 960-, 120-, and
240-byte records accommodate the respective encoded nodes and lookup
metadata. Keeping the types separate avoids padding every leaf to branch
size. Conservatively allow up to twice the compact $6D$ storage budget
for PIR record padding and indexing, where $D=64N$: the scanned databases
together are budgeted at $D_{\mathrm{PIR}}=12D$. This extra factor is a
planning allowance, not a measured packing ratio.

The representative proof visits $r=8+1+1=10$ nodes. In every round, look
up the next node's hash in all three type databases and discard nonmatches.
This keeps the observed database selection independent of the next node's
type, even when that type is not yet known to the client. Each round scans
their combined $12D$ bytes and returns one complete needed node. Dummy
rounds after termination maintain the same round count within this path
budget. This is 30 typed lookups over ten rounds, not
143 independent scans of a piece database. With work linear in scanned
bytes and the baseline batching efficiency retained, the work multiplier
is $f=rD_{\mathrm{PIR}}/D=120$, giving
\[
 m_{\mathrm{InsPIRe}}=
 \left\lceil\frac{Lf}{q}\right\rceil
 =20{,}000\text{--}30{,}000,
 \qquad q=40\text{--}60\text{ account lookups/sec}.
\]
At \$60,268.80 per GPU-year, this is \$1.205B--\$1.808B/year, about
$7.15\times$ below the earlier piece-query estimate ($143\cdot6=858$
times baseline work). Both numerator and baseline $q$ count complete
keyword lookups, so the existing cuckoo candidate queries are not charged
a second time. This is a proposed layout and linear-work extrapolation;
its record packing, lookup indexing, batching, and multi-GPU placement
still require benchmarking. It stores shared trie nodes, rather than
materializing a separate complete proof for every account.

All annual costs use the retained hourly rates and 8,760 hours/year.
Proof egress, updates, preprocessing, failover, and padding beyond the
representative path budget are not priced. The non-PIR rows depend on
the memory-bound assumption above; InsPIRe depends on the proposed
whole-node layout and its compute-scaling assumptions.

\paragraph{TEE-based ORAM.}
In this approach, the TEE protects the computation and its data, while ORAM
hides which records are accessed from an observer outside the enclave. Our
company, Oblivious Labs, Inc., has developed neOMAP, 
the fastest known TEE-backed oblivious
key-value store~\cite{neomap}.
%and our proposed MicrORAM
%work (Thrust~2) aims to further reduce the computational overhead within
%this model.
\begin{pros}
\item \emph{Almost native performance.} TEE-based ORAM~\cite{neomap} performs best among the
  approaches in our comparison, and its performance is comparable to native
  execution: at Ethereum scale, the query latency can be as small as
  30--50\,$\mu$s, and a single machine can achieve 33K--400K queries per
  second, depending on whether queries are batched\footnote{The neOMAP numbers in
  \Cref{tab:comparison} are measured in a setting similar to InsPIRe's
  deployment and do not include fetching Merkle proofs.
The prototype we developed in our earlier grant
  with the Ethereum Foundation additionally supports fetching Merkle proofs.}.
\item \emph{Logarithmic scaling.} 
The performance scales logarithmically w.r.t. 
the database size, and thus the implementation enjoys high performance
both for moderate and large databases. 
%and thus the implementation demonstrates   
%High throughput per machine makes TEE-based
%  ORAM attractive for services that must support many users at low cost.
\end{pros}
\begin{cons}
\item \emph{Trust in hardware.} Privacy depends on the security of the TEE
  and its hardware vendor.
\end{cons}

\paragraph{Private information retrieval.}
PIR allows a client to retrieve a record while keeping its choice of record
hidden from the server. The Ethereum Foundation's GPU-based InsPIRe
implementation~\cite{efinspire} provides a production-quality implementation
with state-of-the-art performance.
\begin{pros}
\item \emph{No trusted hardware.} Privacy does not rely on trusted hardware
or threshold-based trust. 
\item \emph{Practical at moderate scale.} Impressively, InsPIRe shows that
  practicality can be achieved for a moderately sized database containing
  all Ethereum accounts (roughly 200M entries of 64 bytes each).
  Benchmarking numbers report roughly 250\,ms end-to-end latency and
  40--60 queries per second per GPU (\Cref{tab:comparison}), response times
  that are reasonable for many applications.
\end{pros}
\begin{cons}
%\item \emph{Higher latency.} Its response times are substantially higher
%  than the microsecond-scale access times of TEE-based ORAM.
\item \emph{Linear scaling.} The main challenge for scan-based PIR is
  scaling to larger databases. Each query, or batch of queries, requires a
  scan of the entire database wrapped inside cryptography, so throughput
  degrades, and the computing bill grows, linearly with the database size.
  Although latency can potentially enjoy better scaling with a carefully
  parallelized implementation, the throughput (per unit of hardware) and
  the computing bill cannot be parallelized away.

At the scale of the full Ethereum state, the
  database and associated cryptographic state may no longer fit on a single
  GPU, which necessitates additional GPU provisioning.
\end{cons}

\begingroup
\widowpenalty=10000
\paragraph{Garbled RAM.}
Garbled RAM combines cryptographic secure computation with ORAM to hide data
accesses. Imprecisely speaking, one can view Garbled RAM as replacing the TEE
(i.e., a hardware-enforced secure computer) in a TEE+ORAM implementation with
two-party or multi-party computation (i.e., a cryptography-enforced secure
computer). Our research group at Carnegie Mellon University has developed
state-of-the-art constructions and open-source implementations, including
PicoGRAM and Zebra~\cite{picogram,zebra}.
\begin{pros}
\item \emph{No trusted hardware.} Like PIR, Garbled RAM hides data accesses
  without relying on a TEE.
However, deployment requires two or more
  servers and a threshold assumption that limits how many can be corrupted
  or collude.
\item \emph{Competitive latency.} Preliminary estimates in
  \Cref{tab:comparison} suggest access latencies comparable to PIR for a
  moderate, account-sized database.
\item \emph{Logarithmic scaling.} Garbled RAM replaces a linear scan with
  logarithmic per-access overhead. Moving from an account-sized database to
  the full Ethereum state (or larger) therefore incurs a very mild increase
  in per-access cost, provided sufficient memory is available. Garbled RAM
  thus offers a promising route to outperforming scan-based PIR at larger
  database sizes and as Ethereum's state continues to grow.
\end{pros}
\begin{cons}
%\item \emph{Non-collusion assumption.} 
\item \emph{Prototype maturity.} Our current implementations are research
  prototypes that have not yet been fully optimized. This is partly because
  Garbled RAM was previously a theoretical line of research, and only
  recently did our work make it concretely efficient~\cite{picogram,zebra}.
Nonetheless our preliminary performance numbers
are encouraging and
suggest that with further algorithmic and systems-level optimizations,
Garbled RAM might be more desirable than scan-based PIR
for large databases (e.g., at full Ethereum scale).
Moreover, because comparatively little effort has gone into optimizing
Garbled RAM so far, there may still be much room for further improvement.
\end{cons}
\par
\endgroup

%\begin{samepage}

\subsection{Proposed Work}
\label{sec:proposed}

\paragraph{Thrust 1: transitioning Garbled RAM into deployable systems.}
Our current Garbled RAM implementations are research prototypes developed
through basic research at CMU. Their promising latency results provide a
starting point, but substantial algorithmic and systems work remains. In
particular, the current implementation has not been optimized for throughput,
so the baseline throughput in the comparison is essentially the inverse of
access latency. We will develop techniques to increase aggregate throughput,
reduce memory consumption, and reduce the amount of cryptographic computation
per access. We will also explore parallel execution and GPU acceleration,
alongside new algorithmic and cryptographic techniques for faster preprocessing.
Together, these improvements aim to turn our research advances into practical
infrastructure for private blockchain access. The Garbled RAM estimates remain
preliminary, and the proposed performance figures in the comparison are targets
for this work.


\paragraph{Thrust 2: MicrORAM.}
For deployments that accept trusted hardware, we will develop MicrORAM, a
TEE-based oblivious map that aims to bring private accesses to blockchain
data down to a few microseconds. MicrORAM builds on our ongoing research
on vectorized ORAM for commodity x86 processors. Hierarchical ORAM has
traditionally lagged behind tree-based ORAM in performance, because its
periodic hash-table rebuilds are computationally expensive and pause
incoming requests. MicrORAM addresses both problems. It rebuilds tables
with AVX-512 vector instructions, which reduces the computational
complexity of rebuilding by a logarithmic factor relative to oblivious
sorting and yields an 80--240$\times$ concrete speedup, and it runs
rebuilds in the background so that requests no longer pause for large
rebuilds. For large
entries, such as nodes of Ethereum's Merkle Patricia Trie, it pairs this
design with a vectorized tree-based ORAM. As the comparison above shows,
our targets for the account-sized database are under two microseconds of
latency without batching and more than two million queries per second per
CPU with batching, which would reduce latency by more than an order of
magnitude relative to neOMAP. We will turn our research prototype into a
production-quality library, integrate it into our oblivious key-value
store for Ethereum state and Merkle proofs, and evaluate it at the scale of
the full Ethereum state.

%\par
%\end{samepage}

\subsection{Existing Open Source Repositories}

\paragraph{TEE+Oblivious RAM.}
Our neOMAP implementation realizes an oblivious key-value
store, and the open source is available
at 
\url{https://github.com/diffpriv/olabs-oram}.

In our earlier grant 
with the Ethereum Foundation, we have tailored our oblivious key-value
store to fetching Ethereum states along with the Merkle proofs.
The implementation is available at %\elaine{FILL}. 
\url{https://github.com/obliviouslabs/oblivious_node}.


\paragraph{Garbled RAM.}
Our Garbled RAM implementations are available at 
\url{https://github.com/picogramimpl/picogram}
and 
\url{https://github.com/zebragram/zebragram}. Currently, 
they are research 
prototypes developed
in our basic research at CMU. 
The proposed work aims to enable
technology transfer of the basic research.  




\section{Technical Approach}


\subsection{Faster Garbled RAM}
\label{sec:gram}

We build on our open-source Garbled RAM implementations,
PicoGRAM~\cite{picogram} and Zebra~\cite{zebra}. PicoGRAM is a practical
Garbled RAM based on the Decisional Diffie-Hellman assumption, and Zebra is
an arithmetic Garbled RAM for large words based on the Decisional Composite
Residuosity assumption. To turn these research prototypes into a
deployable service for private access to Ethereum data, we will extend
them to multiple parties, make their preprocessing faster, and accelerate
their cryptographic computation with GPUs and AVX-512 vector instructions.

\paragraph{Extending Garbled RAM to multiple parties.}
PicoGRAM and Zebra are two-party Garbled RAM schemes. 
%Because the privacy
%of Garbled RAM rests on a threshold assumption about how many servers may
%collude (\Cref{sec:comparison}), 
For the Ethereum deployment scenario we
will initially aim for a setting with 10--20 parties, which spreads trust
across a larger set of servers. Our recent work~\cite{mparithgarble} has
already developed the cryptographic groundwork for this extension. Its
techniques make multi-party arithmetic garbling practically efficient,
whereas such garbling was previously mostly of theoretical interest. However,
our earlier work~\cite{mparithgarble}
instantiates the new arithmetic garbling techniques only for circuits. The next step
is therefore to combine the cryptographic techniques of our 
recent work~\cite{mparithgarble}
with the RAM techniques of PicoGRAM~\cite{picogram} and Zebra~\cite{zebra}. 
The new, 
efficient
arithmetic garbling will 
allow us to reduce the bandwidth for
moderate to large entries blocks --- by expressing each block 
as one or more arithmetic words, rather than garbling it bit by bit.  
This approach is also suitable for 
Ethereum's setting records of 64 bytes to kilobytes, 
\elaine{double check}
depending on whether Merkle proofs are attached. 

\paragraph{Faster preprocessing.}
Our current Garbled RAM has a preprocessing phase that generates
multiplication triples and garbles the circuitry that is later evaluated
online. Preprocessing can be performed offline and is embarrassingly
parallel, so it does not count toward the online cost of a query. It is
nonetheless worth making preprocessing more practical: in our current
implementation, which has received little optimization so far,
preprocessing requires at least an order of magnitude more work than the
online phase, and this work adds to the computing bill of running the
service.

Our current construction uses pseudorandom correlation generators (PCGs)
based on Learning Parity with Noise (LPN) to produce the authenticated
randomness and multiplication triples needed for preprocessing. A PCG lets
the parties locally expand short correlated seeds into a large amount of
correlated randomness, which minimizes communication. However, the local
expansion of standard LPN-based PCGs is computationally heavy, and it is
the main reason why preprocessing costs so much more than the online
phase. We will replace this component with a PCG based on
Ring-LPN~\cite{ringlpnpcg}, which exploits the structure of polynomial
rings to generate batches of correlations far more efficiently. Prior work
demonstrates an order-of-magnitude speedup, generating roughly 50,000
triples per second per core instead of roughly 7,000, while reducing the
seed material from gigabytes to megabytes~\cite{ringlpnpcg}. We will adapt
Ring-LPN-based PCGs to our framework and benchmark them, with the goal of
removing preprocessing as an end-to-end bottleneck.

\paragraph{GPU and AVX-512 acceleration.}
Both garbling and preprocessing are highly parallel and should therefore
benefit from hardware acceleration. In PicoGRAM, the garbler's most
expensive operations are elliptic-curve scalar multiplications, each
taking roughly 120,000 CPU cycles, and the garbler's computation can be
distributed across many cores~\cite{picogram}. Similarly, a PCG expands a
large batch of correlations at once. We will accelerate garbling and
preprocessing on GPUs, and on CPUs with AVX-512 vector instructions,
building on our experience vectorizing ORAM with AVX-512 in MicrORAM
(Thrust~2).

\paragraph{Prototyping and evaluation.}
Combining the techniques above, we will develop an efficient multi-party
Garbled RAM implementation and release it as open source, as we did for
PicoGRAM and Zebra. We will evaluate it first on the account dataset
(${\sim}206.8$M records). Our current implementation keeps all of the
garbled circuitry in memory~\cite{picogram}, so we will then develop
further memory optimizations and evaluate the implementation on the full
Ethereum state ($2^{32}$ records). For both datasets, we will report a
performance comparison with the Ethereum Foundation's PIR-based approach,
its GPU-based InsPIRe implementation~\cite{efinspire}.

\subsection{MicrORAM: Micro-Second Level Oblivious Accesses for Blockchain Data}
\label{sec:microram}

%A TEE protects the contents of memory, but an attacker outside the enclave
%can still observe which memory locations the enclave accesses, for example
%through page faults and cache side channels~\cite{controlledchannel}. ORAM
%hides these access patterns. Inside a TEE, every step of the ORAM
%algorithm, including the steps that reorganize data, must itself be
%oblivious.

ORAM constructions fall primarily into two categories: hierarchical and
tree-based. Hierarchical ORAM~\cite{GO96,oram11} organizes records into hash
tables of exponentially increasing size. Every access writes the requested
record to the smallest table, so the smaller tables periodically fill up and
must be rebuilt into larger ones. Tree-based ORAM~\cite{pathoramjrnl}
avoids periodic rebuilds altogether. Each access retrieves the record's
location from a \emph{position map} and performs a fixed amount of
maintenance, termed \emph{eviction}. Tree-based ORAM has traditionally been
favored for concrete performance~\cite{neomap}, 
because the rebuilds of hierarchical ORAM
introduce two major drawbacks. The first is high computational overhead:
the state-of-the-art hierarchical ORAM implementation~\cite{h2o2} 
rebuilds its tables with expensive
oblivious sorting, at $O(\log^3 N)$ amortized work per access for $N$
records. The second is high tail latency, because incoming requests must
pause during large rebuilds. 

We observe, however, that these limitations
of hierarchical ORAM 
stem from specific implementation choices rather than from the
hierarchical paradigm itself. We will therefore design novel algorithmic
techniques, including micro-architectural-level optimizations and careful
deamortization, that avoid the parts of hierarchical ORAM that have
limited its performance. Micro-architectural optimizations enable
vectorized rebuilding that  
avoids the high overhead of oblivious sorting, 
and deamortization removes
the pauses for large rebuilds. These techniques will allow hierarchical
ORAM to fully deliver its practical performance, and we expect them to
yield an order-of-magnitude performance improvement over our current
neOMAP~\cite{neomap}. Part of this improvement comes from the fact that
hierarchical ORAM does not need the \emph{recursive} position map of
tree-based ORAM. In a tree-based scheme, reading a record's entry in the
position map directly would reveal which record is accessed, so the
position map is itself stored in a smaller ORAM, whose position map is
stored in a yet smaller ORAM, and so on; these recursive lookups are one
of the performance bottlenecks of tree-based schemes.

\paragraph{Vectorized rebuilding.}
%Although an attacker can observe memory access patterns, data movement
%within the processor's registers remains hidden. 
We rely on the modern processors' 
vector instructions 
(e.g., AVX-512)
to design a rebuild algorithm that reduces the
computational complexity by a logarithmic factor compared with existing
methods which use oblivious
sorting. Instead of sorting the records, MicrORAM routes them to their
destination buckets through a butterfly network with $\Theta(\log n)$
stages on $n$ records, whereas bitonic sort~\cite{bitonicsort} requires
$\Theta(\log^2 n)$ layers. Each bucket is stored in a \emph{byte-sliced}
layout, in which one 512-bit register holds the same byte of all 64 records
in the bucket, so a single permutation instruction moves that byte of every
record to its destination. The secret routing decisions never leave the
registers, and the memory trace depends only on public table sizes.
Together with other optimizations, this design achieves an
80--240$\times$ concrete speedup in rebuilding.

\paragraph{Deamortization: rebuilding in the background.}
To solve the tail-latency problem, 
we design
%extend classic
%deamortization~\cite{OS97} with 
new demortization techniques that execute rebuilds
concurrently on background workers, so that incoming requests no longer
pause for large rebuilds. Each request reads a consistent snapshot of the
published tables, and a completed rebuild atomically replaces its input
tables, so a request sees either the old tables or their complete
replacement.
We expect our new deamortization technique
to achieve a 52\% reduction in memory
footage and 29\% reduction 
in computation overhead in comparison
with known standard techniques~\cite{OS97}.

\paragraph{Further algorithmic optimizations.}
When data entries are small, we expect our hierarchical ORAM alone to
deliver the best performance. With larger entries, however, such as when
Merkle proofs are retrieved along with Ethereum state, 
%and the entries
%include nodes of Ethereum's Merkle Patricia Trie, 
we expect a hybrid
approach to perform better: a tree-based ORAM stores the entries, and our
hierarchical ORAM serves as its position map. Part of the reason is that
a tree-based ORAM performs best when each tree node occupies one
or more full cache lines. Small entries either leave part of each cache
line unused, which wastes memory bandwidth, or require more entries per
tree node to fill it, which adds blocks to every accessed path and thus
adds eviction work. 
Large entries, in contrast, fill whole cache lines
even with few entries per node.
In our hybrid design, the hierarchical ORAM
stores only the small position-map records, which keeps large payloads
out of its hash tables and their rebuilds and replaces the recursive
position map that a tree-based ORAM would otherwise need. We also
vectorize the tree's eviction: with a byte-sliced layout similar to that
of our rebuilds, vector permutation instructions move every block on the
accessed path to its destination, and multiple workers share the eviction
of large entries by each handling a disjoint range of payload bytes. 
%A
%unified interface will select the hierarchical scheme for small entries
%and the hybrid scheme for large ones.

\paragraph{Prototyping and evaluation.}
We currently have a work-in-progress implementation of MicrORAM. Although
it is not yet feature-complete, it already achieves at least an
order-of-magnitude speedup over our current neOMAP~\cite{neomap}. In the
proposed work, we will complete this implementation and deliver a
production-quality library with a common interface for small and large
entries. At Ethereum scale, we target a latency of 2--3\,$\mu$s per query
and a throughput of 1M--2M queries per second per CPU with batching
(\Cref{tab:comparison}). We will integrate it into our oblivious key-value store for
Ethereum state and Merkle proofs, and test it on both the account dataset
(${\sim}206.8$M records) and the full Ethereum state ($2^{32}$ records),
including the retrieval of Merkle paths. This faster version will allow us
to upgrade the oblivious server that we delivered to the Ethereum
Foundation in our earlier grant. 
We will also work on formal verification of the security of our scheme
and implementation.
%We will also strengthen its security argument, including measuring whether
%the AVX-512 instructions we rely on execute in operand-independent time on
%the target processors.

\subsection{Open Source Commitment}

All software artifacts developed by Oblivious Labs, Inc. so far
have been released through open source under MIT license.
We have not filed any patents.
We will continue this commitment for the work supported by this grant. 



